\documentclass[11pt]{article}
\usepackage[a4paper,margin=1in]{geometry}
\usepackage[T1]{fontenc}
\usepackage[utf8]{inputenc}
\usepackage{lmodern}
\usepackage{amsmath,amssymb,amsthm,mathtools}
\usepackage{microtype}
\usepackage{fancyhdr}
\usepackage{array,longtable,enumitem}
\setlength{\parindent}{1.5em}
\setlength{\parskip}{0.18em}
\emergencystretch=4em
\pagestyle{fancy}
\fancyhf{}
\cfoot{\thepage}

\usepackage[ngerman]{babel}
\lhead{Heinrich Weber, Lehrbuch der Algebra, Band III}
\rhead{deutsche Abschrift}
\begin{document}
\begin{center}
{\large Erster Abschnitt.}\\[0.7em]
{\Large Die elliptischen Integrale.}
\end{center}

\section*{\S\,7. Elliptische Raumkurven vierter Ordnung.}

Man kann algebraische Funktionen einer Veränderlichen auch durch Gleichungen zwischen mehreren Veränderlichen darstellen, wenn man die Anzahl der Gleichungen entsprechend vergrößert. Nimmt man z.~B. drei Variable $x,y,z$ und läßt zwischen ihnen zwei Gleichungen
\begin{equation}
\varphi(x,y,z)=0,\qquad \psi(x,y,z)=0
\tag{1}
\end{equation}
bestehen, so kann man aus diesen beiden Gleichungen z.~B. $x$ eliminieren und erhält eine Gleichung zwischen $y$ und $z$, durch die $y$ als algebraische Funktion von $z$ definiert ist. Es kann dann, wenn man gewisse Ausnahmefälle ausschließt, $x$ und jede rationale Funktion von $x,y,z$ rational durch $y$ und $z$ dargestellt werden. Die Integrale der Form
\begin{equation}
\int F(x,y,z)\,dz,
\tag{2}
\end{equation}
in denen $F$ eine rationale Funktion bedeutet, gehören dann zu dem durch (1) dargestellten algebraischen Gebilde. Nimmt man $x,y,z$ als Cartesische Koordinaten im Raume an, so stellt (1) eine Raumkurve als den Durchschnitt zweier Flächen $\varphi=0$, $\psi=0$ dar, und die Integrale (2) gehören zu dieser Raumkurve. Die Kurve heißt wieder elliptisch, wenn diese Integrale elliptisch sind.

Wir wollen diese Betrachtungen auf die Raumkurven vierter Ordnung erster Spezies anwenden, d.~h. auf die Kurven vierter Ordnung, die sich als vollständiger Durchschnitt zweier Flächen zweiten Grades darstellen lassen.

Wir führen wieder homogene Koordinaten $x_1,x_2,x_3,x_4$ ein und nehmen die Gleichungen zweier gegebenen Flächen zweiten Grades in der Form an:
\begin{equation}
\varphi(x_1,x_2,x_3,x_4)=\sum a_{ik}x_i x_k,\qquad
\psi(x_1,x_2,x_3,x_4)=\sum b_{ik}x_i x_k.
\tag{3}
\end{equation}
Die beiden Flächen bestimmen ein Flächenbüschel zweiter Ordnung $\xi\varphi=\eta\psi$. Alle Flächen des Büschels schneiden sich in einer Raumkurve vierter Ordnung, die wir die Grundkurve nennen. Man erhält dieselbe Kurve und dasselbe Büschel, wenn man die Funktionen $\varphi$ und $\psi$ durch $\varphi',\psi'$ ersetzt, die aus $\varphi,\psi$ mittels der linearen Substitution $\begin{pmatrix}m&n\\ p&q\end{pmatrix}$, also durch
\begin{equation}
\varphi'=m\varphi+n\psi,\qquad
\psi'=p\varphi+q\psi,
\tag{4}
\end{equation}
abgeleitet wird, deren Determinante $r=mq-np$ von Null verschieden ist. Dadurch ergibt sich die Identität
\begin{equation}
\xi\varphi+\eta\psi=\xi'\varphi'+\eta'\psi',
\tag{5}
\end{equation}
worin
\begin{equation}
\xi=m\xi'+p\eta',\qquad \eta=n\xi'+q\eta'
\tag{6}
\end{equation}
eine lineare Transformation der Variablen $\xi,\eta$ darstellt. Außer dieser binären Substitution kommt noch eine lineare quaternäre Substitution der Variablen $x_1,x_2,x_3,x_4$ (Koordinatentransformation) in Betracht, deren Determinante wir gleich $1$ annehmen können.

Jede quadratische Form besitzt diesen letzteren Transformationen gegenüber eine Invariante, nämlich die Hessesche Determinante (Bd.~I, \S\S~62, 63, 66), und wir erhalten also als Invariante der Form (5):
\begin{equation}
f(\xi,\eta)=\sum\pm(\xi a_{11}+\eta b_{11})(\xi a_{22}+\eta b_{22})(\xi a_{33}+\eta b_{33})(\xi a_{44}+\eta b_{44}).
\tag{7}
\end{equation}
Dies ist eine binäre biquadratische Form der Variablen $\xi,\eta$, die wir, entwickelt, so darstellen:
\begin{equation}
f(\xi,\eta)=a_0\xi^4+a_1\xi^3\eta+a_2\xi^2\eta^2+a_3\xi\eta^3+a_4\eta^4.
\tag{8}
\end{equation}
Die Koeffizienten dieser biquadratischen Form $a_0,a_1,a_2,a_3,a_4$ ändern sich nicht bei einer Koordinatentransformation. Sie heißen daher simultane Invarianten des Formenpaares $\varphi,\psi$ ($a_0$ und $a_4$ sind die Determinanten von $\varphi$ und von $\psi$).

Die $a_i$ ändern sich, wenn man die $\xi,\eta$ einer Substitution (6) unterwirft. Bildet man aber die Invarianten der Form (8) (nach Algebra, Bd.~I, \S~70), so erhält man Funktionen der Koeffizienten, die auch diesen Substitutionen gegenüber invariant sind, die also nicht zu den individuellen Formen $\varphi,\psi$, sondern zu dem ganzen Büschel und also auch zu ihrem Durchschnitt, der Grundkurve, gehören. Wir nennen sie Invarianten der Grundkurve. In bezug auf das Formensystem $\varphi,\psi$ werden sie auch Kombinanten genannt. Die Grundkurve hat also zwei Invarianten:
\begin{equation}
A=a_2^2-3a_1a_3+12a_0a_4,
\qquad
B=27a_1^2a_4+27a_0a_3^2+2a_2^3-72a_0a_2a_4-9a_1a_2a_3,
\tag{9}
\end{equation}
aus denen man die Diskriminante $D$ nach der Formel
\begin{equation}
27D=4A^3-B^2
\tag{10}
\end{equation}
ableitet. In bezug auf die Koeffizienten von $\varphi$ und $\psi$ sind die Invarianten $A,B,D$ von den Graden $8,12,24$.

Das Formensystem $\varphi,\psi$ hat zwei simultane quadratische Kovarianten, die man ebenso wie die Kovarianten $C$ (in Algebra I, \S~65) ableitet. Die erste von ihnen ist, wenn wir $\psi_i=\frac12\frac{\partial\psi}{\partial x_i}$ setzen:
\begin{equation}
\Phi=
\begin{vmatrix}
a_{11}&a_{12}&a_{13}&a_{14}&\psi_1\\
a_{21}&a_{22}&a_{23}&a_{24}&\psi_2\\
a_{31}&a_{32}&a_{33}&a_{34}&\psi_3\\
a_{41}&a_{42}&a_{43}&a_{44}&\psi_4\\
\psi_1&\psi_2&\psi_3&\psi_4&0
\end{vmatrix},
\tag{11}
\end{equation}
und die zweite $\Psi$ erhält man daraus, indem man $a_{ik}$ mit $b_{ik}$ und $\psi$ mit $\varphi$ vertauscht.

Die Gleichung $\Phi=0$ drückt eine Fläche zweiten Grades aus (die nicht zum Büschel gehört), die der geometrische Ort der Punkte $x$ ist, deren Polaren in bezug auf $\psi$ die Fläche $\varphi$ berühren.

Wir nehmen jetzt an, daß die Diskriminante $D$ von Null verschieden ist. Dann lassen sich die beiden Funktionen $\varphi,\psi$ simultan in die Summe von vier Quadraten transformieren:
\begin{equation}
\varphi=\alpha_1y_1^2+\alpha_2y_2^2+\alpha_3y_3^2+\alpha_4y_4^2,\qquad
\psi=\beta_1y_1^2+\beta_2y_2^2+\beta_3y_3^2+\beta_4y_4^2,
\tag{12}
\end{equation}
worin die $y_i$ lineare Funktionen der $x_i$ sind. Dies ist die kanonische Form des Funktionenpaares. Es wird dann:
\begin{equation}
f(\xi,\eta)=(\xi\alpha_1+\eta\beta_1)(\xi\alpha_2+\eta\beta_2)(\xi\alpha_3+\eta\beta_3)(\xi\alpha_4+\eta\beta_4),
\tag{13}
\end{equation}
und die Funktionen $\Phi,\Psi$ erhalten die kanonische Form:
\begin{equation}
\Phi=\sum_{i=1}^4 \beta_i^2\!\!\prod_{j\ne i}\alpha_j\,y_i^2,\qquad
\Psi=\sum_{i=1}^4 \alpha_i^2\!\!\prod_{j\ne i}\beta_j\,y_i^2.
\tag{14}
\end{equation}
Die Determinante des Systems (12), (14), als lineare Gleichungen für $y_1^2,y_2^2,y_3^2,y_4^2$ betrachtet, ist
\begin{equation}
\begin{vmatrix}
\alpha_1&\alpha_2&\alpha_3&\alpha_4\\
\beta_1&\beta_2&\beta_3&\beta_4\\
\beta_1^2\alpha_2\alpha_3\alpha_4&
\beta_2^2\alpha_1\alpha_3\alpha_4&
\beta_3^2\alpha_1\alpha_2\alpha_4&
\beta_4^2\alpha_1\alpha_2\alpha_3\\
\alpha_1^2\beta_2\beta_3\beta_4&
\alpha_2^2\beta_1\beta_3\beta_4&
\alpha_3^2\beta_1\beta_2\beta_4&
\alpha_4^2\beta_1\beta_2\beta_3
\end{vmatrix}.
\tag{15}
\end{equation}
Sie verschwindet, wenn $\alpha_i:\alpha_k=\beta_i:\beta_k$ wird, weil dann die beiden ersten Kolonnen miteinander proportional werden, und ebenso wenn $\alpha_i:\alpha_k=\beta_i:\beta_k$ wird. Es ist daher (15) als ganze Funktion der $\alpha,\beta$ betrachtet, teilbar durch
\begin{equation}
\alpha_i\beta_k-\alpha_k\beta_i=(\alpha_i\beta_k),
\tag{16}
\end{equation}
und folglich auch durch das Produkt aller dieser Faktoren
\begin{equation}
(\alpha_1\beta_2)(\alpha_1\beta_3)(\alpha_1\beta_4)(\alpha_2\beta_3)(\alpha_2\beta_4)(\alpha_3\beta_4)=\Delta.
\tag{17}
\end{equation}
Aus der Vergleichung der Grade ergibt sich, daß sich die beiden Funktionen (15), (17) nur durch einen Zahlenfaktor unterscheiden können, und dieser ergibt sich aus der Annahme $\beta_1=0$, $\alpha_2=0$, $\alpha_4=0$ gleich $1$. Die Determinante (15) ist also dem Produkt $\Delta$ gleich, und das Quadrat von $\Delta$ ist die Diskriminante $D$ der Funktion $f(\xi,\eta)$, also gleich der Funktion $D$, von der wir angenommen haben, daß sie von Null verschieden sei. Es lassen sich also die $y_1^2,y_2^2,y_3^2,y_4^2$ linear durch $\varphi,\psi,\Phi,\Psi$ darstellen.

Die Kovarianten $\Phi,\Psi$ ändern sich, wenn $\varphi$ und $\psi$ nach (4) durch $\varphi',\psi'$ ersetzt werden. Sie gehören also nicht zu der Grundkurve, sondern zu dem Formenpaar $\varphi,\psi$. Aber es gehen bei dieser Substitution $\Phi$ und $\Psi$ als lineare Funktionen von $y_1^2,y_2^2,y_3^2,y_4^2$ in lineare Funktionen von $\varphi,\psi,\Phi,\Psi$ über. Diese Ausdrücke sind ziemlich kompliziert. Wir brauchen sie hier aber nur für die Punkte der Grundkurve, also für $\varphi=0$, $\psi=0$, und unter dieser Voraussetzung werden sie sehr einfach. Es ist da nämlich
\begin{equation}
\begin{aligned}
\Phi'&=\sum (p\alpha_1+q\beta_1)^2(m\alpha_2+n\beta_2)(m\alpha_3+n\beta_3)(m\alpha_4+n\beta_4)y_1^2,\\
\Psi'&=\sum (m\alpha_1+n\beta_1)^2(p\alpha_2+q\beta_2)(p\alpha_3+q\beta_3)(p\alpha_4+q\beta_4)y_1^2.
\end{aligned}
\tag{18}
\end{equation}
und daraus ist nach (12) und (14) abzuleiten:
\[
\Phi'=M\Phi+N\Psi,\qquad \Psi'=P\Phi+Q\Psi,
\]
worin die Koeffizienten $M,N,P,Q$ noch zu bestimmen sind. Wir können dies durch Rechnung ausführen. Ohne Rechnung ergibt sich das Resultat auf folgendem Wege:

Die $M,N$ sind ganze Funktionen zweiten Grades von $p,q$ und dritten Grades von $m,n$. Nimmt man dann $p:q=m:n$, also $r=mq-np=0$, so gehen $\Phi'$ und $\partial\Phi'/\partial p$ nach (18) in lineare Verbindungen von $\varphi$ und $\psi$ über und verschwinden also. Daraus folgt, daß $M$ und $N$ durch $r^2$ teilbar sind, und die Quotienten sind lineare Funktionen von $m$ und $n$. Da aber für die beiden Fälle
\[
\begin{pmatrix}m&n\\ p&q\end{pmatrix}=
\begin{pmatrix}1&0\\ 0&1\end{pmatrix},
\qquad
\begin{pmatrix}0&1\\ 1&0\end{pmatrix}
\]
$\Phi'$ in $\Phi$ und in $\Psi$ übergehen muß, und da man dieselbe Betrachtung auf $\Psi'$ anwenden kann, so folgt:
\begin{equation}
\Phi'=r^2(m\Phi+n\Psi),\qquad
\Psi'=r^2(p\Phi+q\Psi).
\tag{19}
\end{equation}

Es werden also (von dem Faktor $r^2$ abgesehen) die Funktionen $\Phi,\Psi$ mit den $\varphi,\psi$ kongredient transformiert. (Über die strenge Begründung dieser Schlüsse sehe man Bd.~I, \S~20.)

Aus diesen Ergebnissen können wir auf einfache Weise, immer unter der Voraussetzung $\varphi=0$, $\psi=0$, die $y_i^2$ durch $\Phi$ und $\Psi$ ausdrücken, also die linearen Gleichungen (12), (14) auflösen. Wenn wir nämlich die Formeln (19) auf die Substitution
\[
\begin{pmatrix}m&n\\ p&q\end{pmatrix}
=
\begin{pmatrix}\beta_1&-\alpha_1\\ 0&1\end{pmatrix}
\]
anwenden, so ist $r=\beta_1$, die Funktion $\psi'$ bleibt ungeändert $=\psi$ und $\varphi'$ geht aus $\varphi$ hervor, wenn
\[
\alpha_1,\alpha_2,\alpha_3,\alpha_4
\]
durch
\[
0,\quad (\alpha_2\beta_1),\quad (\alpha_3\beta_1),\quad (\alpha_4\beta_1)
\]
ersetzt werden. Es wird also nach (14)
\[
\Phi'=\beta_1^2(\alpha_2\beta_1)(\alpha_3\beta_1)(\alpha_4\beta_1)y_1^2
\]
und folglich nach (19):
\begin{equation}
\begin{aligned}
(\alpha_2\beta_1)(\alpha_3\beta_1)(\alpha_4\beta_1)y_1^2&=\alpha_1\Phi-\beta_1\Psi,\\
(\alpha_1\beta_2)(\alpha_3\beta_2)(\alpha_4\beta_2)y_2^2&=\alpha_2\Phi-\beta_2\Psi,\\
(\alpha_1\beta_3)(\alpha_2\beta_3)(\alpha_4\beta_3)y_3^2&=\alpha_3\Phi-\beta_3\Psi,\\
(\alpha_1\beta_4)(\alpha_2\beta_4)(\alpha_3\beta_4)y_4^2&=\alpha_4\Phi-\beta_4\Psi.
\end{aligned}
\tag{20}
\end{equation}

Das Produkt des konstanten Faktors auf der linken Seite ist die Diskriminante $D$, und man erhält also nach (13)
\begin{equation}
D y_1^2y_2^2y_3^2y_4^2=f(\Phi,-\Psi)
=a_0\Phi^4-a_1\Phi^3\Psi+a_2\Phi^2\Psi^2-a_3\Phi\Psi^3+a_4\Psi^4,
\tag{21}
\end{equation}
und die $y_i^2$ sind also die linearen Faktoren dieser biquadratischen Form.

Es gibt noch eine weitere simultane Kovariante der Form $\varphi,\psi$, nämlich die Jacobische Funktionaldeterminante der vier Funktionen $\varphi,\psi,\Phi,\Psi$. Wir bezeichnen sie so:
\begin{equation}
K=
\begin{vmatrix}
\varphi_1&\varphi_2&\varphi_3&\varphi_4\\
\psi_1&\psi_2&\psi_3&\psi_4\\
\Phi_1&\Phi_2&\Phi_3&\Phi_4\\
\Psi_1&\Psi_2&\Psi_3&\Psi_4
\end{vmatrix}.
\tag{22}
\end{equation}
Bildet man sie für die kanonische Form, so ergibt sich nach (12), (14), (15)
\[
K=-\Delta y_1y_2y_3y_4,
\]
und da $\Delta^2=D$ ist:
\[
K^2=D y_1^2y_2^2y_3^2y_4^2,
\]
also haben wir nach (21)
\begin{equation}
K^2=f(\Phi,-\Psi).
\tag{23}
\end{equation}

Die biquadratische Form $f(\Phi,-\Psi)$ hat zwei Kovarianten, die wir wie in Bd.~I, \S~70 so bezeichnen:
\begin{equation}
H=\frac13
\begin{vmatrix}
\dfrac{\partial^2f}{\partial\Phi^2}&-\dfrac{\partial^2f}{\partial\Phi\,\partial\Psi}\\[0.5em]
\dfrac{\partial^2f}{\partial\Psi\,\partial\Phi}&-\dfrac{\partial^2f}{\partial\Psi^2}
\end{vmatrix},
\qquad
T=\frac1{12}
\begin{vmatrix}
f'(\Phi)&f'(-\Psi)\\
H'(\Phi)&H'(-\Psi)
\end{vmatrix}.
\tag{24}
\end{equation}
Sie sind vom vierten und sechsten Grade in $\Phi,\Psi$, also vom achten und zwölften Grade in den $x_i$. Es sind Kovarianten der Grundkurve (Kombinanten von $\varphi$ und $\psi$), und es besteht zwischen ihnen die Relation:
\begin{equation}
H^3-48AH f^2-64B f^3=-27T^2.
\tag{25}
\end{equation}

Das allgemeinste zu der Grundkurve gehörige Integral (2) geht durch Einführung homogener Variabler $z=x_3:x_4$ in folgende Form über:
\[
\int F'(x,y,z)\frac{x_3\,dx_4-x_4\,dx_3}{x_4^2},
\]
oder wenn man
\[
F(x_1,x_2,x_3,x_4)=x_4^2\,\frac{F'(x,y,z)}{\varphi_1\psi_2-\varphi_2\psi_1}
\]
setzt, in
\[
\int F(x_1,x_2,x_3,x_4)\,
\frac{x_3\,dx_4-x_4\,dx_3}{\varphi_1\psi_2-\varphi_2\psi_1},
\]
worin $F$ eine homogene Funktion $0$ten Grades ist. Der einfachste Fall ist $F=1$, und wir betrachten also das Integral erster Gattung:
\begin{equation}
u=\int \frac{x_3\,dx_4-x_4\,dx_3}{\varphi_1\psi_2-\varphi_2\psi_1}.
\tag{26}
\end{equation}

Nun bilden wir nach dem Multiplikationssatze der Determinanten das Produkt $K(x_3dx_4-x_4dx_3)$, also wegen $\varphi=0,\psi=0$:
\[
\begin{vmatrix}
\varphi_1&\varphi_2&\varphi_3&\varphi_4\\
\psi_1&\psi_2&\psi_3&\psi_4\\
\Phi_1&\Phi_2&\Phi_3&\Phi_4\\
\Psi_1&\Psi_2&\Psi_3&\Psi_4
\end{vmatrix}
\begin{vmatrix}
1&0&0&0\\
0&1&0&0\\
x_1&x_2&x_3&x_4\\
dx_1&dx_2&dx_3&dx_4
\end{vmatrix}
=
\begin{vmatrix}
\varphi_1&\varphi_2&0&0\\
\psi_1&\psi_2&0&0\\
\Phi_1&\Phi_2&d\Phi&\Phi\\
\Psi_1&\Psi_2&d\Psi&\Psi
\end{vmatrix}
\]
\[
=(\varphi_1\psi_2-\varphi_2\psi_1)(\Phi\,d\Psi-\Psi\,d\Phi),
\]
und folglich ergibt sich aus (26) und (23)
\begin{equation}
u=\int\frac{\Phi\,d\Psi-\Psi\,d\Phi}{K}
=\int\frac{\Phi\,d\Psi-\Psi\,d\Phi}{\sqrt{f(\Phi,-\Psi)}}.
\tag{27}
\end{equation}
Wendet man hierauf die Transformation des \S~5 an, indem man $x,y$ durch $\Phi,-\Psi$ ersetzt, so folgt
\[
u=\sqrt3\int \frac{f\,dH-H\,df}{\sqrt{-(H^3-48Af^2H-64Bf^3)}\,f},
\]
und wenn man in \S~5, (17) $\mu=3$, also
\[
z=-\frac{3H}{4f}
\]
setzt:
\[
u=3\int\frac{dz}{\sqrt{4z^3-g_2z-g_3}},
\qquad
g_2=12\cdot9A,\quad g_3=-4\cdot27B.
\]

\section*{\S\,8. Das Jacobische Transformationsprinzip.}

Die Jacobische Transformationstheorie stellt sich im allgemeinen die Aufgabe, ein elliptisches Differential durch eine algebraische Substitution auf ein anderes elliptisches Differential zurückzuführen. Indem wir hier die Grundlagen dieser Theorie auseinandersetzen, bedienen wir uns der Darstellung des elliptischen Differentials in homogenen Variablen [\S~1, (4)]:
\begin{equation}
\frac{x\,dy-y\,dx}{\sqrt{f(x,y)}},
\tag{1}
\end{equation}
und substituieren darin für $x,y$ zwei teilerfremde ganze homogene Funktionen gleichen, aber beliebigen Grades $n$ zweier neuer Variablen $\xi,\eta$:
\begin{equation}
x=U(\xi,\eta),\qquad y=V(\xi,\eta).
\tag{2}
\end{equation}
Aus den Gleichungen
\[
nU=\xi\frac{\partial U}{\partial\xi}+\eta\frac{\partial U}{\partial\eta},\qquad
dU=d\xi\frac{\partial U}{\partial\xi}+d\eta\frac{\partial U}{\partial\eta},
\]
\[
nV=\xi\frac{\partial V}{\partial\xi}+\eta\frac{\partial V}{\partial\eta},\qquad
dV=d\xi\frac{\partial V}{\partial\xi}+d\eta\frac{\partial V}{\partial\eta},
\]
folgt sodann
\begin{equation}
U\,dV-V\,dU=H(\xi\,d\eta-\eta\,d\xi),
\tag{3}
\end{equation}
wenn $H$ die Funktionaldeterminante
\begin{equation}
H=\frac1n\left(
\frac{\partial U}{\partial\xi}\frac{\partial V}{\partial\eta}
-
\frac{\partial V}{\partial\xi}\frac{\partial U}{\partial\eta}
\right)
\tag{4}
\end{equation}
bedeutet.

Danach geht das elliptische Differential (1) in das folgende über:
\begin{equation}
H\,\frac{\xi\,d\eta-\eta\,d\xi}{\sqrt{f(U,V)}}.
\tag{5}
\end{equation}
Damit nun dieses Differential wieder die Form eines elliptischen erhält, ist erforderlich, daß von der Funktion $f(U,V)$, deren Grad $4n$ ist, sich ein quadratischer Faktor $T^2$ vom Grade $4n-4$ absondern lasse, also, wenn $\varphi(\xi,\eta)$ eine Funktion vierten Grades bedeutet, daß
\begin{equation}
f(U,V)=T^2\varphi(\xi,\eta)
\tag{6}
\end{equation}
werde, wodurch, wenn
\begin{equation}
\frac{T}{H}=M
\tag{7}
\end{equation}
gesetzt wird, das Differential (5) in
\begin{equation}
\frac1M\,\frac{\xi\,d\eta-\eta\,d\xi}{\sqrt{\varphi(\xi,\eta)}}
\tag{8}
\end{equation}
übergeht. Es läßt sich nun nachweisen, daß, sobald die Bedingung (6) erfüllt ist, $H$ durch $T$ teilbar, und also, da der Grad beider Funktionen derselbe ist, $M$ eine Konstante wird. Diese Konstante heißt der Multiplikator der Transformation.

Bemerken wir nämlich, daß, da $f(x,y)$ keinen quadratischen Faktor enthält, die beiden Funktionen
\[
\frac{\partial f}{\partial x},\qquad \frac{\partial f}{\partial y}
\]
keinen gemeinschaftlichen Teiler haben, und daß infolgedessen, weil $U$ und $V$ teilerfremd sind, auch
\[
\frac{\partial f}{\partial U},\qquad \frac{\partial f}{\partial V}
\]
keinen gemeinschaftlichen Teiler (in Beziehung auf $\xi,\eta$) haben, so lassen sich zwei Funktionen $\alpha,\beta$ von $\xi,\eta$ so bestimmen, daß
\[
\alpha\frac{\partial f}{\partial U}+\beta\frac{\partial f}{\partial V}
\]
zu einer beliebig gegebenen Funktion $T$ teilerfremd ist. Man kann z.~B. $\alpha$ teilerfremd zu $T$ und $\beta$ durch die in $\frac{\partial f}{\partial U}$ nicht aufgehenden Teiler von $T$ teilbar, dagegen durch die gemeinschaftlichen Teiler von $T$ und $\frac{\partial f}{\partial U}$ unteilbar annehmen. Wenn wir nun die Gleichung (6), die wir als erfüllt voraussetzen, nach $\xi,\eta$ differentieren, so folgt:
\[
\frac{\partial f}{\partial U}\frac{\partial U}{\partial\xi}
+\frac{\partial f}{\partial V}\frac{\partial V}{\partial\xi}=TX,
\]
\[
\frac{\partial f}{\partial U}\frac{\partial U}{\partial\eta}
+\frac{\partial f}{\partial V}\frac{\partial V}{\partial\eta}=TY,
\]
und daraus durch Auflösung in bezug auf $\frac{\partial f}{\partial U}$, $\frac{\partial f}{\partial V}$:
\[
H\left(\alpha\frac{\partial f}{\partial U}+\beta\frac{\partial f}{\partial V}\right)=TZ,
\]
worin $X,Y,Z$ ganze homogene Funktionen von $\xi,\eta$ sind. Aus der letzten Gleichung aber schließt man, daß $H$ durch $T$ teilbar sein muß.

Demnach ist unser ganzes Problem enthalten in der Gleichung (6), die nichts anderes besagt, als daß die Funktion $4n$ten Grades $f(U,V)$, $2n-2$ quadratische Faktoren enthalten soll. Zur Befriedigung der hieraus folgenden $2n-2$ Bedingungen hat man die $2n+2$ in $U,V$ enthaltenen Koeffizienten zur Verfügung, so daß vier von diesen unbestimmt bleiben. Dies war vorauszusehen, da für $\xi,\eta$ beliebige homogene lineare Funktionen von $\xi,\eta$ eingeführt werden können. Man kann diese vier übrigbleibenden Konstanten dazu verwenden, um das Differential (8) in eine Normalform zu bringen. Da aber die Bedingungsgleichungen für die Koeffizienten von $U,V$ nicht linear sind, so gibt es für einen gegebenen Transformationsgrad mehrere Transformationen.

Wir werden diesem Transformationsproblem später von einer ganz anderen Seite her wieder begegnen und weit tiefer darauf eingehen müssen. Es soll daher auf die Einzelheiten des algebraischen Problems hier nicht näher eingegangen werden; dagegen wollen wir durch zwei Beispiele das Gesagte veranschaulichen.

\section*{\S\,9. Die Transformation zweiten Grades.}

Wir setzen, um die elliptischen Differentiale in der Normalform zu erhalten:
\[
f(x,y)=xy(x-y)(x-\lambda^2y),\qquad
\varphi(\xi,\eta)=\xi\eta(\xi-\eta)(\xi-\varkappa^2\eta),
\]
und es seien $U,V$ vom zweiten Grade.

Die Gleichung (6), \S~8, wird jetzt, da $T$ vom zweiten Grade ist,
\begin{equation}
UV(U-V)(U-\lambda^2V)
=(a\xi+b\eta)^2(a'\xi+b'\eta)^2\xi\eta(\xi-\eta)(\xi-\varkappa^2\eta),
\tag{1}
\end{equation}
und es müssen zwei der vier Faktoren zweiten Grades auf der linken Seite dieser Gleichung Quadrate sein. Die große Zahl der hierin liegenden Möglichkeiten wollen wir dadurch noch beschränken, daß wir voraussetzen, $\eta$ und $y$ sollen gleichzeitig verschwinden, also $V$ durch $\eta$ teilbar sein. Dies gibt (von einem konstanten Faktor abgesehen) für $V$ die folgenden drei Möglichkeiten:
\[
1.\; V=\xi\eta,\qquad
2.\; V=\eta(\xi-\eta),\qquad
3.\; V=\eta(\xi-\varkappa^2\eta).
\]
Jeder dieser drei Fälle umfaßt nun wieder drei Unterfälle, indem von den drei übrigen Faktoren $U$, $U-V$, $U-\lambda^2V$ irgend zwei als Quadrate angenommen werden können. Von diesen letzteren drei Fällen gehen zwei ineinander über durch Vertauschung von $V,\lambda^2$ mit $V:\lambda^2,1:\lambda^2$.

Wir wollen zwei für die Folge besonders wichtige unter diesen Transformationen vollständig durchführen.

\subsection*{1. Die Gausssche Transformation.}

\[
V=\xi\eta,\qquad U=(a\xi+b\eta)^2.
\]
Wenn nun noch $U-\lambda^2V$ ein Quadrat sein soll, so ist
\[
\lambda^2=4ab
\]
zu setzen, und es wird
\[
U-\lambda^2V=(a\xi-b\eta)^2.
\]
Da $U-V$ sodann durch $\xi-\eta$ und durch $\xi-\varkappa^2\eta$ teilbar sein muß, so ergeben sich, wenn man $\xi=\eta$, $\xi=\varkappa^2\eta$ setzt, aus $U=V$ die Bedingungen
\[
a+b=\pm1,\qquad a\varkappa^2+b=\pm\varkappa,
\]
woraus, wenn die oberen Zeichen genommen werden,
\[
a=\frac1{1+\varkappa},\qquad
b=\frac{\varkappa}{1+\varkappa},\qquad
\lambda=\frac{2\sqrt{\varkappa}}{1+\varkappa},
\]
\[
V=\xi\eta,\qquad
U=\left(\frac{\xi+\varkappa\eta}{1+\varkappa}\right)^2,\qquad
U-\lambda^2V=\left(\frac{\xi-\varkappa\eta}{1+\varkappa}\right)^2,
\]
\[
U-V=\frac{(\xi-\eta)(\xi-\varkappa^2\eta)}{(1+\varkappa)^2},
\]
\[
T=\frac{\xi^2-\varkappa^2\eta^2}{(1+\varkappa)^3},\qquad
H=\frac{\xi^2-\varkappa^2\eta^2}{(1+\varkappa)^2},\qquad
M=\frac1{1+\varkappa}.
\]
Setzt man also wieder
\[
\frac{y}{x}=z,\qquad \frac{\eta}{\xi}=\zeta,
\]
so haben wir das Resultat, daß durch die Substitution
\begin{equation}
z=\frac{(1+\varkappa)^2\zeta}{(1+\varkappa\zeta)^2},
\qquad
\lambda^2=\frac{4\varkappa}{(1+\varkappa)^2}
\tag{2}
\end{equation}
die Transformation
\begin{equation}
\frac{dz}{\sqrt{z(1-z)(1-\lambda^2z)}}
=
\frac{(1+\varkappa)d\zeta}{\sqrt{\zeta(1-\zeta)(1-\varkappa^2\zeta)}}
\tag{3}
\end{equation}
geleistet wird.

\subsection*{2. Die Landensche Transformation.}

\[
V=a\eta(\xi-\eta),\qquad U=\xi(\xi-\varkappa^2\eta).
\]
Die Bedingungen, daß $U-V$, $U-\lambda^2V$ Quadrate sind, lauten
\[
4a=(\varkappa^2+a)^2,\qquad 4a\lambda^2=(\varkappa^2+a\lambda^2)^2,
\]
woraus
\[
\lambda(\varkappa^2+a)=\varkappa^2+a\lambda^2,
\]
oder durch $\lambda-1$ dividiert
\[
a\lambda=\varkappa^2;
\]
dies in eine der obigen Gleichungen eingesetzt, gibt
\[
4\lambda=\varkappa^2(1+\lambda)^2,
\]
und durch Auflösung dieser quadratischen Gleichung, wenn $\varkappa'=\sqrt{1-\varkappa^2}$ gesetzt ist:
\[
\lambda=\frac{1-\varkappa'}{1+\varkappa'},\qquad
a=(1+\varkappa')^2,
\]
\[
U-V=\{\xi-(1+\varkappa')\eta\}^2,\qquad
U-\lambda^2V=\{\xi-(1-\varkappa')\eta\}^2,
\]
\[
H=(1+\varkappa')^2\{\xi-(1+\varkappa')\eta\}\{\xi-(1-\varkappa')\eta\},
\]
\[
T=(1+\varkappa')\{(\xi-\eta)^2-\varkappa'^2\eta^2\},\qquad
M=\frac1{1+\varkappa'},
\]
also durch die Substitution
\begin{equation}
z=\frac{(1+\varkappa')^2\zeta(1-\zeta)}{1-\varkappa^2\zeta},
\qquad
\lambda=\frac{1-\varkappa'}{1+\varkappa'},
\tag{4}
\end{equation}
die Umformung:
\begin{equation}
\frac{dz}{\sqrt{z(1-z)(1-\lambda^2z)}}
=
\frac{(1+\varkappa')\,d\zeta}{\sqrt{\zeta(1-\zeta)(1-\varkappa^2\zeta)}}.
\tag{5}
\end{equation}
Wendet man auf die linke Seite dieser Gleichung wieder die Gausssche Transformation an, indem man in den Formeln (2), (3) $\xi,\lambda$ durch $z,\lambda$ und $z$ durch eine Variable $\eta$ ersetzt, so ist, wie aus (2) und (4) folgt, $\lambda$ in (3) durch $\varkappa$ zu ersetzen, und durch Kombination von (5) und (3) findet man:
\begin{equation}
\frac{d\eta}{\sqrt{\eta(1-\eta)(1-\varkappa^2\eta)}}
=
\frac{2\,d\zeta}{\sqrt{\zeta(1-\zeta)(1-\varkappa^2\zeta)}},
\tag{6}
\end{equation}
und die Verbindung von (2) und (4) ergibt zwischen den Variablen $\eta$ und $\zeta$ den Zusammenhang:
\begin{equation}
\eta=\frac{4\zeta(1-\zeta)(1-\varkappa^2\zeta)}
{(1-\varkappa^2\zeta^2)^2}.
\tag{7}
\end{equation}
Die Kombination der beiden Transformationen zweiten Grades gibt also die Multiplikation mit $2$.

\section*{\S\,10. Die Transformation dritten Grades.}

Als zweites Beispiel betrachten wir die Transformation dritten Grades. Die Gleichung
\begin{equation}
UV(U-V)(U-\lambda^2V)=T^2\xi\eta(\xi-\eta)(\xi-\varkappa^2\eta)
\tag{1}
\end{equation}
fordert, daß jeder der vier Faktoren dritten Grades $U,V,U-V,U-\lambda^2V$ durch einen der Linearfaktoren $\xi,\eta,\xi-\eta,\xi-\varkappa^2\eta$ teilbar und daß der Quotient ein Quadrat sei. Wir setzen also
\begin{equation}
U=\xi(a\xi+b\eta)^2,\qquad
V=\eta(a'\xi+b'\eta)^2,
\tag{2}
\end{equation}
und verlangen noch, daß
\begin{equation}
\frac{U-V}{\xi-\eta},\qquad
\frac{U-\lambda^2V}{\xi-\varkappa^2\eta}
\tag{3}
\end{equation}
die Quadrate linearer Funktionen werden.

Von den beiden letzten Forderungen folgt die eine aus der anderen, wenn wir $U,V$ so einrichten, daß sie, von konstanten Faktoren abgesehen, $U$ und $V$ ineinander übergehen durch die Vertauschung von $\xi,\eta$ mit $\varkappa^2\eta,\xi$. Durch diese Vertauschung geht aber
\[
\xi(a\xi+b\eta)^2,\qquad \eta(a'\xi+b'\eta)^2
\]
über in
\[
\varkappa^2\eta(b\xi+a\varkappa^2\eta)^2,\qquad
\xi(b'\xi+a'\varkappa^2\eta)^2,
\]
und unsere Forderung ist erfüllt, wenn
\begin{equation}
\varkappa^2=\frac{bb'}{aa'}
\tag{4}
\end{equation}
ist. Wenn dann
\begin{equation}
\lambda^2=\varkappa^2\left(\frac{ab}{a'b'}\right)^2
\tag{5}
\end{equation}
ist, so geht durch diese Vertauschung
\begin{equation}
\frac{U-V}{\xi-\eta}\quad\hbox{in}\quad
\frac{b'^2}{a^2}\frac{U-\lambda^2V}{\xi-\varkappa^2\eta}
\tag{6}
\end{equation}
über. Nachdem dies festgesetzt, ist nur noch die Bedingung zu erfüllen, daß die erste der beiden Größen (3) das Quadrat einer linearen Funktion wird, die wir mit $(\alpha^2\xi-\beta^2\eta)$ bezeichnen, also:
\begin{equation}
U-V=(\xi-\eta)(\alpha^2\xi-\beta^2\eta)^2
\tag{7}
\end{equation}
oder
\[
U-\xi(\alpha^2\xi-\beta^2\eta)^2
=
V-\eta(\alpha^2\xi-\beta^2\eta)^2,
\]
und da $U$ durch $\xi$, $V$ durch $\eta$ teilbar ist, so muß diese Funktion durch $\xi\eta$ teilbar sein; setzen wir sie $=\xi\eta(m\xi+n\eta)$, so ergibt sich aus (2)
\[
(a\xi+b\eta)^2=(\alpha^2\xi-\beta^2\eta)^2+\eta(m\xi+n\eta),
\]
\[
(a'\xi+b'\eta)^2=(\alpha^2\xi-\beta^2\eta)^2+\xi(m\xi+n\eta).
\]
Die Vergleichung der Koeffizienten ergibt
\[
a=\alpha^2,\qquad b=-\beta^2+\frac{m}{2\alpha^2},\qquad b^2=\beta^4+n,
\]
\[
b'=\beta^2,\qquad a'=-\alpha^2+\frac{n}{2\beta^2},\qquad a'^2=\alpha^4+m,
\]
und aus den beiden letzten Gleichungen jeder Reihe:
\[
m^2=4\alpha^2(n\alpha^2+m\beta^2),\qquad
n^2=4\beta^2(n\alpha^2+m\beta^2).
\]
Wenn man also
\[
m=h\alpha,\qquad n=h\beta
\]
setzt, so wird
\[
h=4\alpha\beta(\alpha+\beta).
\]
Hiernach lassen sich $a,b,a',b'$ durch die beiden Größen $\alpha,\beta$ ausdrücken in der Weise:
\begin{equation}
a=\alpha^2,\qquad b=\beta^2+2\alpha\beta,\qquad
a'=\alpha^2+2\alpha\beta,\qquad b'=\beta^2,
\tag{8}
\end{equation}
und danach aus (4), (5)
\begin{equation}
\varkappa^2=\frac{\beta^3}{\alpha^3}\frac{\beta+2\alpha}{\alpha+2\beta},
\qquad
\lambda^2=\frac{\alpha}{\beta}\frac{\beta+2\alpha}{\alpha+2\beta},
\tag{9}
\end{equation}
oder durch Multiplikation und Division dieser beiden Gleichungen:
\begin{equation}
\sqrt{\lambda\varkappa}=\frac{\beta}{\alpha}\frac{\beta+2\alpha}{\alpha+2\beta},
\qquad
\frac{\varkappa^3}{\lambda}=\frac{\beta^4}{\alpha^4}.
\tag{10}
\end{equation}
Durch Elimination von $\beta:\alpha$ erhält man eine Gleichung zwischen $\varkappa,\lambda$, die man die Modulargleichung nennt, deren Grad die Anzahl der verschiedenen Transformationen dritten Grades angibt. Sie nimmt die einfachste Gestalt an, wenn man setzt:
\begin{equation}
\sqrt[4]{\varkappa}=u,\qquad \sqrt[4]{\lambda}=v.
\tag{11}
\end{equation}
Dann werden die Gleichungen (10)
\[
u^2v^2=\frac{\beta}{\alpha}\frac{\beta+2\alpha}{\alpha+2\beta},
\qquad
\frac{u^3}{v}=\frac{\beta}{\alpha},
\]
also durch Einsetzen des Wertes von $\beta:\alpha$ in die erste Gleichung und Beseitigung des Faktors $u^2$:
\begin{equation}
v^4-u^4+2u^3v^3-2uv=0,
\tag{12}
\end{equation}
eine Gleichung vom vierten Grade.

Man erhält ferner [ohne weitläufige Rechnung durch Vergleichung der Koeffizienten der höchsten Potenz von $\xi$ in (1)]
\[
T=\frac{a}{b'}(a\xi+b\eta)(a'\xi+b'\eta)(\alpha^2\xi-\beta^2\eta)(\beta^2\xi-\varkappa^2\alpha^2\eta),
\]
\[
H=\frac13\left(
\frac{\partial U}{\partial\xi}\frac{\partial V}{\partial\eta}
-\frac{\partial V}{\partial\xi}\frac{\partial U}{\partial\eta}
\right)=\frac{a'}a\,T,
\]
woraus man leicht nach (8) findet:
\begin{equation}
M=\frac{a}{a'}=\frac{\alpha}{\alpha+2\beta}
=\frac{v}{v+2u^3}.
\tag{13}
\end{equation}
Setzt man den hieraus sich ergebenden Ausdruck
\[
v=\frac{2u^3M}{1-M}
\]
in die Gleichung (12) ein, so ergibt sich für $M$ eine Gleichung vierten Grades, die Multiplikatorgleichung, die die Modulargleichung ersetzen kann. Man erhält aber diese Gleichung einfacher auf folgendem Wege. Nach (13) ist
\[
\frac1M-1=2\frac{\beta}{\alpha},\qquad
\frac1M+3=\frac{2(\beta+2\alpha)}{\alpha},
\]
also nach (9):
\[
\left(\frac1M-1\right)^3\left(\frac1M+3\right)
=16\varkappa^2\frac{\alpha+2\beta}{\alpha},
\]
und folglich nach (13)
\[
\frac{16\varkappa^2}{M}
=\left(\frac1M-1\right)^3\left(\frac1M+3\right),
\]
oder geordnet
\begin{equation}
\frac1{M^4}-\frac6{M^2}+\frac{8(1-2\varkappa^2)}{M}-3=0.
\tag{14}
\end{equation}

Drücken wir also unsere Formeln durch $u,v$ aus, so ist das Ergebnis dieser Betrachtung das folgende: Durch die Substitution
\[
z=\frac{\zeta\bigl[(v^2+2vu^3)+u^6\zeta\bigr]^2}
{\bigl[v^2+(u^6+2vu^3)\zeta\bigr]^2}
\]
wird die Transformation bewirkt
\[
\frac{dz}{\sqrt{z(1-z)(1-v^8z)}}
=
\frac{v+2u^3}{v}\,
\frac{d\zeta}{\sqrt{\zeta(1-\zeta)(1-u^8\zeta)}},
\]
falls zwischen $u,v$ die Modulargleichung (12) besteht.

Zu einem gegebenen $u$ ergibt die Modulargleichung vier Werte von $v$, also vier verschiedene Transformationen dritten Grades.

\section*{\S\,11. Die drei Gattungen elliptischer Integrale.}

Es sei jetzt
\begin{equation}
f(x)=a_0x^4+a_1x^3+a_2x^2+a_3x+a_4
\tag{1}
\end{equation}
eine ganze Funktion dritten oder vierten Grades, und
\begin{equation}
\Omega(x)=\int\frac{\Phi(x)\,dx}{\sqrt{f(x)}}=\int R(x)\,dx,
\tag{2}
\end{equation}
wenn zur Abkürzung
\begin{equation}
R(x)=\frac{\Phi(x)}{\sqrt{f(x)}}
\tag{3}
\end{equation}
gesetzt ist, ein elliptisches Integral. Es bedeutet darin $x$ eine unbeschränkt veränderliche (auch komplexe) Größe, und $\sqrt{f(x)}$ hat für jeden Wert von $x$, für den $f(x)$ nicht verschwindet, zwei entgegengesetzte Werte. Wir verstehen unter einem Punkt nicht einen Wert von $x$ allein, sondern ein zusammengehöriges Wertepaar von $x$ und $\sqrt{f(x)}$, also einen Wert von $x$ mit einem bestimmten Vorzeichen von $\sqrt{f(x)}$. Ist daher $x_0$ irgend ein Wert von $x$, so gibt es zwei Punkte $x_0$, wenn $f(x_0)$ von Null verschieden ist, aber nur einen, wenn $f(x_0)=0$ ist. Für $x=\infty$ bestimmen wir die Punkte nach dem Vorzeichen von $\sqrt{f(x)}:x^2$, und wir haben also zwei Punkte $x=\infty$, wenn $f(x)$ vom vierten, und nur einen, wenn $f(x)$ vom dritten Grade ist.

Nach bekannten Sätzen der Funktionentheorie gibt es, wenn $f(x_0)$ von Null verschieden ist, eine in einem gewissen Bereich konvergente Entwickelung nach dem Taylorschen Lehrsatz:
\begin{equation}
R(x)=A_m(x-x_0)^m+A_{m+1}(x-x_0)^{m+1}+A_{m+2}(x-x_0)^{m+2}+\cdots,
\tag{4}
\end{equation}
worin $m$ eine positive oder negative ganze Zahl oder auch Null ist, während $A_m,A_{m+1},A_{m+2},\ldots$ Konstanten bedeuten, von denen $A_m$ nicht verschwindet. Dies ist eine Entwickelung nach steigenden Potenzen. Ergänzend tritt hinzu, falls $f(x)$ vom vierten (nicht vom dritten) Grade ist, die Entwickelung nach fallenden Potenzen
\begin{equation}
R(x)=C_mx^{-m-2}+C_{m+1}x^{-m-3}+C_{m+2}x^{-m-4}+\cdots,
\tag{5}
\end{equation}
von der wir sagen, daß sie in der Umgebung eines unendlich fernen Punktes gilt.

Wenn aber $f(x_0)$ verschwindet, so erhält die Entwickelung von $R(x)$ die Form:
\begin{equation}
R(x)=A_m(x-x_0)^{m-\frac12}+A_{m+1}(x-x_0)^{m-\frac12+1}
+A_{m+2}(x-x_0)^{m-\frac12+2}+\cdots,
\tag{6}
\end{equation}
und wenn $f(x)$ vom dritten Grade ist, so gilt in der Umgebung des unendlich fernen Punktes die Entwickelung
\begin{equation}
R(x)=C_mx^{-m-\frac32}+C_{m+1}x^{-m-\frac52}
+C_{m+2}x^{-m-\frac72}+\cdots.
\tag{7}
\end{equation}

Aus den Entwickelungen (4) bis (7) können wir die Entwickelungen von $\Omega$ ableiten, wobei eine additive Konstante unbestimmt bleibt.

Die Entwickelung von $\Omega$ ist dann ebenfalls eine Potenzreihe, wozu, wenn $m$ negativ ist, in den Fällen (4) oder (5) noch ein Glied $A_{-1}\log(x-x_0)$, $C_{-1}\log x$ tritt.

Wir nennen nun $x_0$ einen Punkt erster Gattung von $\Omega$, wenn $m$ in diesen Entwickelungen nicht negativ ist. Dann ist $\Omega$ im Punkte $x_0$ endlich (auch im Falle $x_0=\infty$).

Der Punkt $x_0$ heißt ein Punkt zweiter Gattung von $\Omega$, wenn in den Entwickelungen (4), (5) $m$ negativ und $A_{-1}=0$, $C_{-1}=0$ ist, und in dem Falle der Entwickelung (6), (7) mit negativem $m$.

Endlich heißt $x_0$ ein Punkt dritter Gattung, wenn in den Entwickelungen (4) oder (5) $A_{-1}$ oder $C_{-1}$ von Null verschieden ist, wenn also in der Entwickelung von $\Omega$ ein logarithmisches Glied vorkommt.

Die Nullpunkte von $f(x)$ können bei dem Integral (2) nicht von der dritten Gattung sein; sie können aber von der dritten Gattung werden, wenn wir das elliptische Integral in der allgemeinen Form
\[
\int\Phi(x,\sqrt{f(x)})\,dx
\]
betrachten.

Das Integral $\Omega$ heißt ein Integral erster Gattung, wenn es nur Punkte erster Gattung hat. Ein solches Integral hat dann für alle endlichen und unendlichen Werte von $x$ einen endlichen Wert.

$\Omega$ heißt ein Integral zweiter Gattung, wenn es nur Punkte erster und zweiter Gattung hat, und $\Omega$ heißt ein Integral dritter Gattung, wenn es neben Punkten erster und zweiter Gattung auch Punkte dritter Gattung hat.

\section*{\S\,12. Darstellung der elliptischen Integrale durch die einfachsten Grundintegrale.}

Um das Integral $\Omega(x)$ durch Integrale von möglichst einfachem Typus darzustellen, zerlegen wir die rationale Funktion $\Phi(x)$ in eine ganze Funktion und in eine Summe von Partialbrüchen, d.~h. wir stellen $\Phi(x)$ dar als eine Summe von Ausdrücken der Form
\[
x^n,\qquad \frac1{(x-\alpha)^n}
\]
mit konstanten Koeffizienten, worin $n$ eine ganze positive Zahl (auch Null) und $\alpha$ einen konstanten Wert bedeutet, für den $\Phi(x)$ unendlich wird.

Setzen wir dann für positive und für negative oder verschwindende $n$
\begin{equation}
S_n(\alpha)=\int\frac{(x-\alpha)^n\,dx}{\sqrt{f(x)}},
\tag{1}
\end{equation}
so wird $\Omega(x)$ eine Summe der Form
\[
\Omega(x)=\sum M S_n(\alpha),
\]
wo sich die Summe auf mehrere verschiedene Werte von $n$ und von $\alpha$ erstrecken kann und die $M$ Konstanten sind.

Die Funktionen $S_n$ lassen sich durch Vermittelung algebraischer Funktionen auf eine geringe Anzahl zurückführen. Zu dem Ende bilden wir
\[
d\!\left((x-\alpha)^n\sqrt{f(x)}\right)
=
\frac{n(x-\alpha)^{n-1}f(x)+(x-\alpha)^n\frac12 f'(x)}
{\sqrt{f(x)}}\,dx,
\]
und wenn wir darin setzen
\[
f(x)=f(\alpha)+(x-\alpha)f'(\alpha)+\frac{(x-\alpha)^2}{2}f''(\alpha)
+\frac{(x-\alpha)^3}{6}f'''(\alpha)
+\frac{(x-\alpha)^4}{24}f''''(\alpha),
\]
\[
f'(x)=f'(\alpha)+(x-\alpha)f''(\alpha)
+\frac{(x-\alpha)^2}{2}f'''(\alpha)
+\frac{(x-\alpha)^3}{6}f''''(\alpha),
\]
so folgt
\begin{equation}
\begin{aligned}
(x-\alpha)^n\sqrt{f(x)}
={}& n f(\alpha)S_{n-1}
+\left(n+\frac12\right)f'(\alpha)S_n\\
&+\frac{n+1}{2}f''(\alpha)S_{n+1}
+\left(\frac n6+\frac14\right)f'''(\alpha)S_{n+2}\\
&+\left(\frac n{24}+\frac1{12}\right)f''''(\alpha)S_{n+3}.
\end{aligned}
\tag{2}
\end{equation}

Wir unterscheiden vier Fälle:

1. $f'''(\alpha)$ und $f(\alpha)$ nicht $=0$ [d.~h. $f(x)$ vom vierten Grade und $\alpha$ keine Wurzel von $f(x)$]. In diesem Falle läßt sich durch die Formel (2) für $n=0$ $S_3$ durch $S_0,S_1,S_2$ ausdrücken, und daher $S_n$ für jedes positive $n$ durch $S_0,S_1,S_2$. Setzt man aber $n=-1,-2,\ldots$, so erhält man $S_{-2},S_{-3},\ldots$ ausgedrückt durch $S_{-1},S_0,S_1,S_2$. Also bleiben in diesem Falle die vier Grundintegrale
\[
S_{-1},\quad S_0,\quad S_1,\quad S_2.
\]

2. $f'''(\alpha)$ nicht $=0$, $f(\alpha)=0$. In diesem Falle kann man noch $S_{-1}$ durch $S_0,S_1,S_2$ ausdrücken und erhält als Grundintegrale
\[
S_0,\quad S_1,\quad S_2.
\]

3. $f'''(\alpha)=0$, $f(\alpha)$ nicht $=0$. Hier sind die Grundintegrale
\[
S_{-1},\quad S_0,\quad S_1.
\]

4. $f'''(\alpha)=0$, $f(\alpha)=0$. Hier sind die Grundintegrale
\[
S_0,\quad S_1.
\]

Das Resultat hiervon ist also:

I. Ist $f(x)$ vom vierten Grade ($a_0=1$), so lassen sich alle Integrale $\Omega(x)$ ausdrücken durch Integrale der Form
\[
S_0=\int\frac{dx}{\sqrt{f(x)}},\qquad
S_1=\int\frac{x\,dx}{\sqrt{f(x)}},\qquad
S_2=\int\frac{x^2\,dx}{\sqrt{f(x)}},
\]
\[
S_{-1}=\int\frac{dx}{(x-\alpha)\sqrt{f(x)}},
\]
worin $S_{-1}$ noch von dem Parameter $\alpha$ abhängig ist.

II. Ist aber $f(x)$ vom dritten Grade, so genügen die drei Integrale
\[
S_0=\int\frac{dx}{\sqrt{f(x)}},\qquad
S_1=\int\frac{x\,dx}{\sqrt{f(x)}},\qquad
S_{-1}=\int\frac{dx}{(x-\alpha)\sqrt{f(x)}}.
\]

In beiden Fällen ist $S_0$ von der ersten Gattung. Ist $f(x)$ vom dritten Grade, so ist $S_1$ von der zweiten und $S_{-1}$ von der dritten Gattung. Ist dagegen $f(x)$ vom vierten Grade, so sind $S_1,S_2$ und $S_{-1}$ von der dritten Gattung. Man kann aber aus $S_1$ und $S_2$ ein Integral zweiter Gattung zusammensetzen. Denn es ist
\[
\frac1{\sqrt{f(x)}}=(x^4+a_1x^3+a_2x^2+a_3x+a_4)^{-1/2}
=x^{-2}-\frac{a_1}{2}x^{-3}+\cdots,
\]
und in der Entwickelung von
\[
\left(x^2+\frac{a_1}{2}x\right):\sqrt{f(x)}
\]
nach fallenden Potenzen kommt kein Glied mit $x^{-1}$ vor. Folglich ist
\[
S_2+\frac{a_1}{2}S_1
\]
ein Integral zweiter Gattung.

Durch Vermittelung einer logarithmischen Funktion kann man aber auch noch die vier Integrale des Falles I auf drei reduzieren. Zu dem Ende bringe man $f(x)$ auf die Form
\[
x^4+a_1x^3+a_2x^2+a_3x+a_4=P^2+ax+b,
\]
worin
\[
P=x^2+px+q
\]
eine quadratische Funktion und $a$ und $b$ Konstanten sind. Es ergibt sich:
\[
p=\frac{a_1}{2},\qquad q=\frac{a_2}{2}-\frac{a_1^2}{8},
\]
\[
a=a_3-\frac{a_1a_2}{2}+\frac{a_1^3}{8},\qquad
b=a_4-\left(\frac{a_2}{2}-\frac{a_1^2}{8}\right)^2.
\]
Nun ist
\begin{equation}
\begin{aligned}
d\log\frac{\sqrt f-P}{\sqrt f+P}
&=
\frac{f'(x)P-2f(x)P'(x)}{f-P^2}\,
\frac{dx}{\sqrt{f(x)}},\\
f-P^2&=ax+b,\qquad f'(x)=2PP'+a.
\end{aligned}
\tag{3}
\end{equation}
und folglich
\[
f'(x)P-2f(x)P'(x)=a(P-2xP')-2bP'=Q
\]
eine quadratische Funktion von $x$, und demnach ergibt sich aus (3)
\begin{equation}
\log\frac{\sqrt f-P}{\sqrt f+P}
=
\int\frac{Q}{ax+b}\frac{dx}{\sqrt{f(x)}}.
\tag{4}
\end{equation}
Die rechte Seite ist aber eine lineare Funktion von $S_0,S_1$ und $S_{-1}$ (für $\alpha=-b/a$). In dem besonderen Falle, wo $a=0$ ist, wird $Q$ vom ersten Grade, und man kann eine lineare Verbindung von $S_0$ und $S_{-1}$ durch eine logarithmische Funktion ausdrücken.

Nehmen wir $a_1$ und $a_3=0$ an, worauf wir den allgemeinen Fall nach \S~4 durch lineare Transformation zurückführen können, so erhalten wir als die Grundintegrale
\[
S_0=\int\frac{dx}{\sqrt{f(x)}},\qquad
S_2=\int\frac{x^2\,dx}{\sqrt{f(x)}},\qquad
S_{-1}=\int\frac{dx}{(x-\alpha)\sqrt{f(x)}},
\]
während $S_1=\int\frac{x\,dx}{\sqrt{f(x)}}$ durch die Substitution $x^2=y$ auf das nicht elliptische Integral
\[
\frac12\int\frac{dy}{\sqrt{y^2+a_2y+a_4}}
\]
zurückgeführt wird. Hier ist $S_0$ von der ersten, $S_2$ von der zweiten, $S_{-1}$ von der dritten Gattung. Für die Legendresche Normalform setzen wir $f(x)=(1-x^2)(1-\varkappa^2x^2)$ und nehmen als Normalintegral der zweiten Gattung nicht $S_2$ selbst, sondern $S_0-\varkappa^2S_2$. Dadurch ergeben sich die Normalintegrale der drei Gattungen:
\[
\int\frac{dx}{\sqrt{(1-x^2)(1-\varkappa^2x^2)}},\qquad
\int\sqrt{\frac{1-\varkappa^2x^2}{1-x^2}}\,dx,\qquad
\int\frac{dx}{(x-\alpha)\sqrt{(1-x^2)(1-\varkappa^2x^2)}},
\]
und wenn man $x=\sin\varphi$ setzt:
\[
\int\frac{d\varphi}{\sqrt{1-\varkappa^2\sin^2\varphi}},\qquad
\int\sqrt{1-\varkappa^2\sin^2\varphi}\,d\varphi,\qquad
\int\frac{d\varphi}{(\sin\varphi-\alpha)\sqrt{1-\varkappa^2\sin^2\varphi}}.
\]
Dasselbe erhält man, wenn man das elliptische Differential in der Normalform
\[
\frac{dz}{\sqrt{z(1-z)(1-\varkappa^2z)}}
\]
annimmt und nach dem Falle II verfährt.

\section*{\S\,13. Das Additionstheorem.}

Das von Euler entdeckte Additionstheorem der elliptischen Integrale besteht in dem Satze, der für die ganze weitere Theorie von fundamentaler Bedeutung ist, daß, wenn von den drei Wertepaaren
\[
x_1,\sqrt{f(x_1)};\qquad x_2,\sqrt{f(x_2)};\qquad x_3,\sqrt{f(x_3)}
\]
zwei als gegeben vorausgesetzt werden, man auf algebraischem Wege das dritte so bestimmen kann, daß die Differentialgleichung
\begin{equation}
\frac{dx_1}{\sqrt{f(x_1)}}+
\frac{dx_2}{\sqrt{f(x_2)}}+
\frac{dx_3}{\sqrt{f(x_3)}}=0
\tag{1}
\end{equation}
befriedigt ist. Darin ist $f(x)$ eine gegebene Funktion vom dritten oder vierten Grade. Es ist dies ein spezieller Fall des großen Abelschen Theorems und soll auch in dieser Weise hier aufgefaßt und abgeleitet werden\footnote{Dieses weitumfassende Theorem findet sich in großartiger Einfachheit abgeleitet in einem kaum zwei Seiten umfassenden Aufsatze im vierten Bande von Crelles Journal: „Démonstration d'une propriété générale d'une certaine classe de fonctions transcendantes“; Oeuvres complètes de N.~H.~Abel. Nouvelle édition, T.~I, p.~515. Ausführliche Darstellung und Anwendungen in der nachgelassenen großen Abhandlung: „Mémoire sur une propriété générale d'une classe très étendue de fonctions transcendantes“.}. Dem Beweise schicken wir folgenden elementaren algebraischen Satz voraus:

Ist $F(x)$ eine ganze rationale Funktion $n$ten Grades ohne mehrfache Faktoren, $F'(x)$ ihre Derivierte, sind ferner $x_1,x_2,\ldots,x_n$ die Wurzeln der Gleichung $F(x)=0$ und $\varphi(x)$ eine ganze Funktion von $x$, deren Grad höchstens $=n-2$, so ist
\begin{equation}
\frac{\varphi(x_1)}{F'(x_1)}
+\frac{\varphi(x_2)}{F'(x_2)}
+\cdots+
\frac{\varphi(x_n)}{F'(x_n)}=0.
\tag{2}
\end{equation}
Der Beweis dieses Satzes ergibt sich unmittelbar aus der Zerlegung des rationalen Bruches
\[
\frac{x\varphi(x)}{F(x)}
\]
in Partialbrüche
\[
\frac{x_1\varphi(x_1)}{F'(x_1)(x-x_1)}
+\frac{x_2\varphi(x_2)}{F'(x_2)(x-x_2)}
+\cdots+
\frac{x_n\varphi(x_n)}{F'(x_n)(x-x_n)},
\]
wenn darin $x=0$ gesetzt wird [Bd.~I, \S~15 (9)].

Es sollen nun $P,Q$ ganze Funktionen von $x$ von den Graden $m$ und $m-2$ bedeuten, und wir fragen nach den Werten von $x,\sqrt{f(x)}$, für die die Funktion
\begin{equation}
P+Q\sqrt{f(x)}
\tag{3}
\end{equation}
verschwindet. Solcher Wertepaare (Punkte) gibt es $2m$, und zwar findet man die Werte von $x$ als Wurzeln der Gleichung $2m$ten Grades:
\begin{equation}
F(x)=P^2-Q^2f(x)=0,
\tag{4}
\end{equation}
und die zugehörigen Werte von $\sqrt{f(x)}$ aus
\begin{equation}
P+Q\sqrt{f(x)}=0.
\tag{5}
\end{equation}

Wir nehmen nun an, die Koeffizienten in den Funktionen $P,Q$ seien veränderlich, entweder unabhängige Veränderliche oder in irgend einer Weise von anderen Veränderlichen abhängig. Es werden dann auch die durch (4), (5) bestimmten Werte $x$ und $\sqrt{f(x)}$ Funktionen dieser Veränderlichen, und (5) läßt sich differentiieren.

Bezeichnen wir mit $\delta$ die Differentiation nach den in den Koeffizienten von $P,Q$ vorkommenden Veränderlichen, wobei $x$ als konstant gilt, mit $dx$ das entsprechende Differential von $x$, wenn $x$ durch (4) oder (5) als Funktion dieser Koeffizienten bestimmt ist, so ergibt die Differentiation von (4)
\[
2P\delta P-2Q\delta Q f(x)+F'(x)\,dx=0,
\]
woraus mit Benutzung von (5):
\begin{equation}
\frac{Q\delta P-P\delta Q}{F'(x)}
=
\frac{dx}{2\sqrt{f(x)}},
\tag{6}
\end{equation}
und da der Zähler auf der linken Seite vom Grade $2m-2$, $F(x)$ vom Grade $2m$ ist, so läßt sich die Formel (2) anwenden, und es folgt
\begin{equation}
\sum \frac{dx}{\sqrt{f(x)}}=0,
\tag{7}
\end{equation}
wenn die Summe auf sämtliche Wurzeln der Gleichung (5) erstreckt ist.

Wir wenden dieses Theorem auf den einfachsten Fall, nämlich $m=2$, an und erhalten dann folgenden Satz:

Wenn
\begin{equation}
x_1,\sqrt{f(x_1)};\quad x_2,\sqrt{f(x_2)};\quad
x_3,\sqrt{f(x_3)};\quad x_4,\sqrt{f(x_4)}
\tag{8}
\end{equation}
vier Wertepaare sind, für die irgend eine Funktion der Form
\begin{equation}
a+bx+cx^2+\sqrt{f(x)}
\tag{9}
\end{equation}
verschwindet, so sind diese Wertepaare nicht gänzlich voneinander unabhängig, sondern es besteht zwischen ihnen die Differentialgleichung
\begin{equation}
\frac{dx_1}{\sqrt{f(x_1)}}+
\frac{dx_2}{\sqrt{f(x_2)}}+
\frac{dx_3}{\sqrt{f(x_3)}}+
\frac{dx_4}{\sqrt{f(x_4)}}=0.
\tag{10}
\end{equation}

Die Abhängigkeit dieser vier Wertepaare, also eine Integration der Differentialgleichung (10), kann aber auch in algebraischer Weise ausgedrückt werden, indem man die Funktion (9) für $x=x_1,x_2,x_3,x_4$ gleich Null setzt und dann $a,b,c$ eliminiert. Man erhält so die Determinantengleichung:
\begin{equation}
\begin{vmatrix}
1&x_1&x_1^2&\sqrt{f(x_1)}\\
1&x_2&x_2^2&\sqrt{f(x_2)}\\
1&x_3&x_3^2&\sqrt{f(x_3)}\\
1&x_4&x_4^2&\sqrt{f(x_4)}
\end{vmatrix}=0.
\tag{11}
\end{equation}
Aus dieser Gleichung kann man $x_1$ und $\sqrt{f(x_1)}$ rational durch $x_2,\sqrt{f(x_2)};\ x_3,\sqrt{f(x_3)};\ x_4,\sqrt{f(x_4)}$ ausdrücken; denn $x_1$ ist die Wurzel einer biquadratischen Gleichung, deren drei andere Wurzeln $x_2,x_3,x_4$ sind, und deren Koeffizienten rational von $x_2,\sqrt{f(x_2)};\ x_3,\sqrt{f(x_3)};\ x_4,\sqrt{f(x_4)}$ abhängen, und $\sqrt{f(x_1)}$ ist mittels (11) rational durch $x_1$ darstellbar. Setzt man $x_4$ einer beliebigen Konstanten gleich, so geht die Differentialgleichung (10) in (1) über und (11) ergibt ihr Integral.

Wir wenden dies Theorem zunächst auf die Normalform an und setzen
\[
f(x)=x(1-x)(1-\varkappa^2x).
\]
Nehmen wir in (11) $x_4=0$ an, so reduziert sich diese Gleichung durch Division mit $\sqrt{x_1x_2x_3}$ auf:
\begin{equation}
\begin{vmatrix}
\sqrt{x_1}&x_1\sqrt{x_1}&\sqrt{(1-x_1)(1-\varkappa^2x_1)}\\
\sqrt{x_2}&x_2\sqrt{x_2}&\sqrt{(1-x_2)(1-\varkappa^2x_2)}\\
\sqrt{x_3}&x_3\sqrt{x_3}&\sqrt{(1-x_3)(1-\varkappa^2x_3)}
\end{vmatrix}=0,
\tag{12}
\end{equation}
oder wenn wir die Determinante nach den Elementen der ersten Zeile entwickeln:
\begin{equation}
\sqrt{x_1}(a+bx_1)+c\sqrt{(1-x_1)(1-\varkappa^2x_1)}=0,
\tag{13}
\end{equation}
worin zur Abkürzung gesetzt ist:
\begin{equation}
\begin{aligned}
a={}&x_2\sqrt{x_2}\sqrt{(1-x_3)(1-\varkappa^2x_3)}
-x_3\sqrt{x_3}\sqrt{(1-x_2)(1-\varkappa^2x_2)},\\
b={}&-\sqrt{x_2}\sqrt{(1-x_3)(1-\varkappa^2x_3)}
+\sqrt{x_3}\sqrt{(1-x_2)(1-\varkappa^2x_2)},\\
c={}&-(x_2-x_3)\sqrt{x_2}\sqrt{x_3}.
\end{aligned}
\tag{14}
\end{equation}
Wenn wir (13) rational machen, so erhalten wir die kubische Gleichung:
\[
x(a+bx)^2-c^2(1-x)(1-\varkappa^2x)=0,
\]
deren drei Wurzeln $x_1,x_2,x_3$ sind. Es ist demnach identisch
\begin{equation}
b^2(x-x_1)(x-x_2)(x-x_3)
=
x(a+bx)^2-c^2(1-x)(1-\varkappa^2x),
\tag{15}
\end{equation}
und indem wir in dieser identischen Gleichung $x=0,1,1:\varkappa^2$ setzen, so folgt:
\begin{equation}
\begin{aligned}
\sqrt{x_1x_2x_3}&=-\frac cb,\\
\sqrt{(1-x_1)(1-x_2)(1-x_3)}&=\frac{b+a}{b},\\
\sqrt{(1-\varkappa^2x_1)(1-\varkappa^2x_2)(1-\varkappa^2x_3)}&=\frac{b+\varkappa^2a}{b}.
\end{aligned}
\tag{16}
\end{equation}
Zur Bestimmung der Vorzeichen in diesen Gleichungen erhalten wir noch aus (13), wenn wir $x_1$ durch $x_2$ und $x_3$ ersetzen und die Ergebnisse multiplizieren:
\[
x_1x_2x_3(a+bx_1)(a+bx_2)(a+bx_3)
=-c^3\sqrt{f(x_1)}\sqrt{f(x_2)}\sqrt{f(x_3)},
\]
und (15) ergibt für $x=-a:b$
\[
(a+bx_1)(a+bx_2)(a+bx_3)
=\frac{c^2}{b^2}(b+a)(b+\varkappa^2a).
\]
Benutzt man noch das Quadrat der ersten Gleichung (16), so folgt hieraus
\[
\sqrt{f(x_1)}\sqrt{f(x_2)}\sqrt{f(x_3)}
=
-\frac{c}{b^3}(b+a)(b+\varkappa^2a).
\]
Dasselbe Resultat ergibt aber die Multiplikation der drei Gleichungen (16), und die Vorzeichen sind also in diesen Formeln so gewählt, daß das Produkt der linken Seiten den durch die Differentialgleichung (1) geforderten Wert von $\sqrt{f(x_1)}$ ergibt. Eine weitere Bestimmung der Vorzeichen ist in (16) der Natur der Sache nach nicht möglich, und diese Gleichungen dienen zur eindeutigen Definition von $\sqrt{x_1}$, $\sqrt{1-x_1}$, $\sqrt{1-\varkappa^2x_1}$, wenn die $\sqrt{x_2}$, $\sqrt{1-x_2}$, $\sqrt{1-\varkappa^2x_2}$, $\sqrt{x_3}$, $\sqrt{1-x_3}$, $\sqrt{1-\varkappa^2x_3}$ als gegeben vorausgesetzt werden.

Um die Gleichungen (16) in die gebräuchliche Form zu bringen, machen wir zunächst in (14) den Zähler von $b$ rational und erhalten
\[
b=-\,\frac{(x_2-x_3)(1-\varkappa^2x_2x_3)}
{\sqrt{x_2}\sqrt{(1-x_3)(1-\varkappa^2x_3)}+
\sqrt{x_3}\sqrt{(1-x_2)(1-\varkappa^2x_2)}},
\]
ferner nach (14):
\[
\frac{b+a}{b}
=\frac{\sqrt{1-x_2}\sqrt{1-x_3}
-\sqrt{x_2x_3}\sqrt{(1-\varkappa^2x_2)(1-\varkappa^2x_3)}}
{1-\varkappa^2x_2x_3},
\]
\[
\frac{b+\varkappa^2a}{b}
=\frac{\sqrt{1-\varkappa^2x_2}\sqrt{1-\varkappa^2x_3}
-\varkappa^2\sqrt{x_2x_3}\sqrt{(1-x_2)(1-x_3)}}
{1-\varkappa^2x_2x_3},
\]
wodurch die Gleichungen (16) leicht in die Form gebracht werden:
\begin{equation}
\begin{aligned}
\sqrt{x_1}&=
-\frac{\sqrt{x_2}\sqrt{(1-x_3)(1-\varkappa^2x_3)}
+\sqrt{x_3}\sqrt{(1-x_2)(1-\varkappa^2x_2)}}
{1-\varkappa^2x_2x_3},\\
\sqrt{1-x_1}&=
\frac{\sqrt{1-x_2}\sqrt{1-x_3}
-\sqrt{x_2x_3}\sqrt{(1-\varkappa^2x_2)(1-\varkappa^2x_3)}}
{1-\varkappa^2x_2x_3},\\
\sqrt{1-\varkappa^2x_1}&=
\frac{\sqrt{1-\varkappa^2x_2}\sqrt{1-\varkappa^2x_3}
-\varkappa^2\sqrt{x_2x_3}\sqrt{(1-x_2)(1-x_3)}}
{1-\varkappa^2x_2x_3}.
\end{aligned}
\tag{17}
\end{equation}
und unsere Betrachtung lehrt, daß, wenn $x_1$ und die Wurzeln $\sqrt{x_1},\sqrt{1-x_1},\sqrt{1-\varkappa^2x_1}$ durch diese Gleichungen als Funktionen von $x_2,x_3$ bestimmt werden, die Differentialgleichung (1) identisch befriedigt ist.

Um auch für die Weierstrasssche Normalform das Additionstheorem in möglichst einfacher Gestalt zu erhalten, setze man
\[
f(x)=4x^3-g_2x-g_3,
\]
und lasse in der Gleichung (11) (nach Division mit $x_4^2$) $x_4$ unendlich groß werden. Es ergibt sich alsdann das Integral der Differentialgleichung (1)
\[
\frac{dx_1}{\sqrt{f(x_1)}}+
\frac{dx_2}{\sqrt{f(x_2)}}+
\frac{dx_3}{\sqrt{f(x_3)}}=0
\]
in der Form
\begin{equation}
\begin{vmatrix}
1&x_1&\sqrt{f(x_1)}\\
1&x_2&\sqrt{f(x_2)}\\
1&x_3&\sqrt{f(x_3)}
\end{vmatrix}=0,
\tag{18}
\end{equation}
oder
\begin{equation}
a+bx_1+c\sqrt{f(x_1)}=0,
\tag{19}
\end{equation}
wenn
\begin{equation}
a=x_2\sqrt{f(x_3)}-x_3\sqrt{f(x_2)},\qquad
b=\sqrt{f(x_2)}-\sqrt{f(x_3)},\qquad
c=x_3-x_2.
\tag{20}
\end{equation}
Es werden dann $x_1,x_2,x_3$ die Wurzeln der kubischen Gleichung
\[
(a+bx)^2-c^2f(x)=0,
\]
und wenn man hierin den Koeffizienten von $x^2$, geteilt durch den von $x^3$, gleich der negativen Summe der Wurzeln setzt, so erhält man
\begin{equation}
x_1+x_2+x_3=\frac14\left(\frac{\sqrt{f(x_2)}-\sqrt{f(x_3)}}{x_2-x_3}\right)^2,
\tag{21}
\end{equation}
und aus (19)
\begin{equation}
\sqrt{f(x_1)}
=
\frac14\left(\frac{\sqrt{f(x_2)}-\sqrt{f(x_3)}}{x_2-x_3}\right)^3
-(x_2+x_3)\frac{\sqrt{f(x_2)}-\sqrt{f(x_3)}}{x_2-x_3}
-\frac{x_3\sqrt{f(x_2)}-x_2\sqrt{f(x_3)}}{x_2-x_3}.
\tag{22}
\end{equation}

\section*{\S\,14. Ursprung der elliptischen Funktionen.}

Wenn in dem elliptischen Differential erster Gattung in der Legendreschen Normalform
\[
\frac12\frac{dz}{\sqrt{z(1-z)(1-\varkappa^2z)}}
\]
die Substitution
\begin{equation}
z=\sin^2\varphi
\tag{1}
\end{equation}
gemacht und sodann die Integration von $\varphi=0$ an ausgeführt wird, so entsteht das elliptische Integral erster Gattung
\begin{equation}
\int_0^\varphi\frac{d\varphi}{\sqrt{1-\varkappa^2\sin^2\varphi}}=u.
\tag{2}
\end{equation}
Jacobi nennt die obere Grenze $\varphi$ dieses Integrals seine Amplitude und schreibt
\begin{equation}
\varphi=\operatorname{am}u.
\tag{3}
\end{equation}
Die trigonometrischen Funktionen dieses Bogens
\begin{equation}
\sin\varphi,\quad \cos\varphi,\quad
\sqrt{1-\varkappa^2\sin^2\varphi}=\Delta\varphi
\tag{4}
\end{equation}
sind es nun, die, als Funktionen von $u$ betrachtet, elliptische Funktionen genannt und von Jacobi mit
\[
\sin\operatorname{am}u,\qquad \cos\operatorname{am}u,\qquad \Delta\operatorname{am}u,
\]
oder, nach Gudermann, kürzer mit
\begin{equation}
\operatorname{sn}u,\qquad \operatorname{cn}u,\qquad \operatorname{dn}u
\tag{5}
\end{equation}
bezeichnet werden. Diese Funktionen hängen von den zwei Argumenten $u$ und $\varkappa^2$ ab, und wenn eine genauere Bezeichnung notwendig ist, wird dafür auch
\[
\operatorname{sn}(u,\varkappa^2),\qquad
\operatorname{cn}(u,\varkappa^2),\qquad
\operatorname{dn}(u,\varkappa^2)
\]
gesetzt.

Wenn in der Formel (1) $\varphi$ in $-\varphi$ verwandelt wird, so geht $u$ in $-u$ über, es ist also $\operatorname{am}(-u)=-\operatorname{am}u$ und folglich
\[
\operatorname{sn}(-u)=-\operatorname{sn}u,\qquad
\operatorname{cn}(-u)=\operatorname{cn}u,\qquad
\operatorname{dn}(-u)=\operatorname{dn}u,
\]
d.~h. es ist $\operatorname{sn}u$ eine ungerade, $\operatorname{cn}u$, $\operatorname{dn}u$ sind gerade Funktionen.

Setzt man also
\[
\frac{d\varphi_2}{\Delta\varphi_2}=du,\qquad
\frac{d\varphi_3}{\Delta\varphi_3}=dv,\qquad
\frac{d\varphi_1}{\Delta\varphi_1}=-d(u+v),
\]
so lassen sich die Formeln (17), \S~13 anwenden, und man erhält die Additionstheoreme:
\begin{equation}
\begin{aligned}
\operatorname{sn}(u+v)&=
\frac{\operatorname{sn}u\,\operatorname{cn}v\,\operatorname{dn}v+
\operatorname{sn}v\,\operatorname{cn}u\,\operatorname{dn}u}
{1-\varkappa^2\operatorname{sn}^2u\,\operatorname{sn}^2v},\\
\operatorname{cn}(u+v)&=
\frac{\operatorname{cn}u\,\operatorname{cn}v-
\operatorname{sn}u\,\operatorname{sn}v\,\operatorname{dn}u\,\operatorname{dn}v}
{1-\varkappa^2\operatorname{sn}^2u\,\operatorname{sn}^2v},\\
\operatorname{dn}(u+v)&=
\frac{\operatorname{dn}u\,\operatorname{dn}v-
\varkappa^2\operatorname{sn}u\,\operatorname{sn}v\,\operatorname{cn}u\,\operatorname{cn}v}
{1-\varkappa^2\operatorname{sn}^2u\,\operatorname{sn}^2v}.
\end{aligned}
\tag{6}
\end{equation}
Man kann daraus die drei elliptischen Funktionen für beliebige Werte des Arguments eindeutig berechnen, wenn sie für irgend ein noch so kleines Gebiet um den Nullpunkt bestimmt sind. Diesen Satz nimmt Weierstrass zum Ausgangspunkt der Theorie der elliptischen Funktionen.

Setzt man
\[
\int_0^{\pi/2}\frac{d\varphi}{\Delta\varphi}=K,
\]
so ist $\operatorname{sn}K=1$, $\operatorname{cn}K=0$, $\operatorname{dn}K=\sqrt{1-\varkappa^2}=\varkappa'$, und es ergeben also die Formeln (6), indem man $v=K$ setzt:
\begin{equation}
\operatorname{sn}(u+K)=\frac{\operatorname{cn}u}{\operatorname{dn}u},\qquad
\operatorname{cn}(u+K)=\frac{-\varkappa'\operatorname{sn}u}{\operatorname{dn}u},\qquad
\operatorname{dn}(u+K)=\frac{\varkappa'}{\operatorname{dn}u},
\tag{7}
\end{equation}
und wenn man diese Formel noch einmal anwendet:
\begin{equation}
\operatorname{sn}(u+2K)=-\operatorname{sn}u,\qquad
\operatorname{cn}(u+2K)=-\operatorname{cn}u,\qquad
\operatorname{dn}(u+2K)=\operatorname{dn}u,
\tag{8}
\end{equation}
und abermals angewandt:
\begin{equation}
\operatorname{sn}(u+4K)=\operatorname{sn}u,\qquad
\operatorname{cn}(u+4K)=\operatorname{cn}u,\qquad
\operatorname{dn}(u+4K)=\operatorname{dn}u.
\tag{9}
\end{equation}
Es haben also die elliptischen Funktionen die Eigenschaft der Periodizität mit den trigonometrischen Funktionen gemein. Sie haben die Periode $4K$ ($\operatorname{dn}u$ auch die Periode $2K$).

In \S~3 haben wir als Beispiel die Transformation, die Formel (1), abgeleitet:
\begin{equation}
\frac{dz}{\sqrt{z(1-z)(1-\varkappa^2z)}}
=
\frac{dx}{i\sqrt{x(1-x)(1-\varkappa'^2x)}},
\tag{10}
\end{equation}
wenn
\begin{equation}
z=\frac{-x}{1-x}
\tag{11}
\end{equation}
war. Setzt man also
\[
\frac{dx}{\sqrt{x(1-x)(1-\varkappa'^2x)}}=i\,du,
\]
so ist, wenn $x$ und $u$ zugleich verschwinden, $x=\operatorname{sn}^2(iu,\varkappa')$, $1-x=\operatorname{cn}^2(iu,\varkappa')$, und aus (10) folgt:
\[
\frac{dz}{\sqrt{z(1-z)(1-\varkappa^2z)}}=du,
\]
also
\[
z=\operatorname{sn}^2(u,\varkappa).
\]
Demnach ergibt (11)
\begin{equation}
\operatorname{sn}(u,\varkappa)=
\frac{-i\,\operatorname{sn}(iu,\varkappa')}{\operatorname{cn}(iu,\varkappa')}
\tag{12}
\end{equation}
(worin das Vorzeichen aus dem speziellen Wert $u=0$ zu bestimmen ist).

Setzt man nun
\[
\int_0^{\pi/2}\frac{d\varphi}{\sqrt{1-\varkappa'^2\sin^2\varphi}}=K',
\]
so ergibt sich aus (8)
\[
\operatorname{sn}(iu+2K',\varkappa')=-\operatorname{sn}(iu,\varkappa'),\qquad
\operatorname{cn}(iu+2K',\varkappa')=-\operatorname{cn}(iu,\varkappa'),
\]
und folglich aus (12)
\[
\operatorname{sn}(u-2iK',\varkappa)=\operatorname{sn}(u,\varkappa)
\]
oder auch, indem man $u$ in $u+2iK'$ verwandelt und die Bezeichnung $\varkappa$ wieder wegläßt:
\begin{equation}
\operatorname{sn}(u+2iK')=\operatorname{sn}u.
\tag{13}
\end{equation}

Es hat also die Funktion $\operatorname{sn}u$ eine doppelte Periodizität. Die eine dieser Perioden, $4K$, ist reell, die andere, $2iK'$, rein imaginär (wenigstens wenn der Modul $\varkappa$ ein positiver echter Bruch ist). Diese doppelte Periodizität ist eine Eigenschaft, die in dem Gebiete der elementaren Funktionen nirgends vorkommt, und so am deutlichsten zeigte, daß man es hier mit einer neuen Funktionenart zu tun hat. Die aus der doppelten Periodizität abgeleiteten Folgerungen sind es zugleich, die, wie die Erfahrung gezeigt hat, weitaus am schnellsten mit befriedigender Strenge in das Innere der Theorie einführen, so daß hier der zweckmäßigste Ausgangspunkt für eine methodische Entwickelung zu suchen ist. Unsere nächsten Betrachtungen werden daher den doppelt periodischen Funktionen im allgemeinen gewidmet sein.


\end{document}
